% Method section of the solver paper. Moved verbatim from draft/bie-siam
% (Section 3, commit cc23eea) on 2026-11-07; to be rewritten (see README.md).
\section{Discretization and candidate computation}\label{sec:discretization}

Candidates are computed on the center cell, as in the transparent boundary
conditions of Fliss \cite[Thm.~4.5]{Fliss2013}: a guided mode restricted to
$\mathcal C_0$ has port Cauchy data of finite-energy half-guide solutions. The
Dirichlet-to-Neumann form of this condition fails at countably many Dirichlet
eigenvalues in the gap \cite[Thm.~4.1]{Fliss2013}, which Robin-to-Robin
conditions avoid \cite{FlissKlindworthSchmidt2015}; we impose instead that the
port states lie in the decaying subspaces themselves. Lead cells are
discretized by boundary integral equations with port Rayleigh amplitudes, the
decaying half-guide solutions by an ordered QZ decomposition, and the center
cell by layer potentials plus background Rayleigh waves. The output is a
candidate $\widehat k$ with coefficients. We make no claim on its accuracy here:
the enclosure paper~\companion bounds the spectral distance of whatever field these
coefficients define. Fix $\beta$ and a wavenumber $k>0$ with $k^2\in G_\beta$.

\subsection{Port Fourier basis and Rayleigh waves}
Let $\beta_m=\beta+2\pi m/d$, $\psi_m(y)=d^{-1/2}\ee^{\ii\beta_my}$, truncated to
$|m|\le M$, $K=2M+1$. With $\kappa_{\mathrm e}=k\rho_{\mathrm e}^{1/2}$ let
$\gamma_m=(\kappa_{\mathrm e}^2-\beta_m^2)^{1/2}$ with $\Im\gamma_m\ge0$, and
$\boldsymbol\gamma=\diag(\gamma_m)$. We assume that $k$ is not a Wood point,
that is, $\gamma_m(k)\ne0$ for all $m\in\ZZ$. Only the candidate search uses
this; the bound of the enclosure paper~\companion does not.

In a vertical background strip, a $\beta$-quasi-periodic Helmholtz solution is
$u=\sum_m(a_mr_m^{x_p}+b_ml_m^{x_p})$ with
$r_m^{x_p}=\ee^{\ii\gamma_m(x-x_p)}\psi_m$ and
$l_m^{x_p}=\ee^{-\ii\gamma_m(x-x_p)}\psi_m$; $(a,b)$ are the right- and
left-going amplitudes anchored at $x_p$. Moving the anchor by $\ell$ multiplies
amplitudes by $\mathsf E(\ell)=\diag(\ee^{\ii\gamma_m\ell})$. For a port with
normal $\pm e_x$, amplitudes map to the Fourier coefficients of (value, normal
derivative) by $\mathsf C_\pm=\left[\begin{smallmatrix}I&I\\
\pm\ii\boldsymbol\gamma&\mp\ii\boldsymbol\gamma\end{smallmatrix}\right]$,
invertible away from Wood points.

\subsection{Layer potentials on a rod}\label{sec:layer}
Let $G^{(\kappa)}$ be the free-space Green function $\frac{\ii}{4}H^{(1)}_0(\kappa|
\mathbf z-\mathbf z'|)$ and $G^{(\kappa)}_\QP$ the $\ee^{\ii\beta d}$-quasi-periodic
one. We use the near-image split
\begin{equation}\label{eq:near-image}
  G^{(\kappa)}_\QP(\mathbf z,\mathbf z')=\sum_{n=-1}^{1}\ee^{\ii\beta nd}G^{(\kappa)}(\mathbf z,
  \mathbf z'+nd\,e_y)+G^{(\kappa)}_{\mathrm{far}}(\mathbf z-\mathbf z'),
\end{equation}
where $G^{(\kappa)}_{\mathrm{far}}$ is a smooth Helmholtz solution away from the
images $|n|\ge2$, since $(\Delta+\kappa^2)G^{(\kappa)}_\QP(\cdot,\mathbf 0)=
-\sum_n\ee^{\ii\beta nd}\delta_{nd\,e_y}$ \cite{Linton1998}. For a curve strictly inside a cell,
a quasi-periodic layer potential is thus three free-space layer potentials plus
a smooth correction near the closed cell, the structure used in
\cite{BarnettGreengard2010,BarnettGreengard2011,ChoBarnett2015}; the smooth
correction and the background are represented together by proxy point sources
outside the cell, with coefficients fitted to the quasi-periodicity on the top
and bottom seams. Since the enclosure paper~\companion checks only the jumps of an explicit
field, the densities may be computed with any adequate approximation of
$G_\QP$; the field evaluated in the enclosure paper~\companion uses the split
\eqref{eq:near-image} and fixed proxies, and all its mismatch enters its own
jumps.

For a rod $\Omega$ with $C^2$ boundary, outward normal $\nu$, and interior
wavenumber $\kappa_{\mathrm c}=k\rho_{\mathrm c}^{1/2}$, we represent the field in
the cell by
\begin{equation}\label{eq:cell-rep}
  u=\mathcal S^{(\kappa_{\mathrm e})}_\QP[\omega]+\mathcal D^{(\kappa_{\mathrm e})}_\QP[\tau]
  +u_{\mathrm b}\ \text{outside }\Omega,\qquad
  u=\mathcal S^{(\kappa_{\mathrm c})}[\omega]+\mathcal D^{(\kappa_{\mathrm c})}[\tau]
  \ \text{in }\Omega,
\end{equation}
with a finite background Rayleigh field $u_{\mathrm b}$. By the jump relations
\cite{ColtonKress1983,Kress1991}, the outside-minus-inside jumps of value and
normal derivative are, with $\eta=(\tau;-\omega)$,
\begin{equation}\label{eq:jump-operator}
  \mathscr J_k[\eta;u_{\mathrm b}]=
  \begin{bmatrix}
    I+D^{(\kappa_{\mathrm e})}_\QP-D^{(\kappa_{\mathrm c})}&-S^{(\kappa_{\mathrm e})}_\QP+S^{(\kappa_{\mathrm c})}\\
    T^{(\kappa_{\mathrm e})}_\QP-T^{(\kappa_{\mathrm c})}&I-D'^{(\kappa_{\mathrm e})}_\QP+D'^{(\kappa_{\mathrm c})}
  \end{bmatrix}\eta
  +\begin{bmatrix}u_{\mathrm b}\\ \partial_\nu u_{\mathrm b}\end{bmatrix}_{\partial\Omega}.
\end{equation}
The transmission conditions are $\mathscr J_k=0$, a second-kind system in which
the hypersingular parts cancel \cite{DominguezLyonTurc2016}. We discretize it by
Nystr\"om quadrature with Kress's logarithmic splitting on a uniform parameter
grid \cite{Kress1991,HaoBarnettMartinssonYoung2014}; $n$ is the number of density
unknowns. By the spectral representation of $G_\QP$, the potentials of a curve
lying left (right) of a line $x=x_p$ are purely right- (left-) going there;
their amplitudes are integrals of $\omega$ and $\tau$ against explicit
exponentials, and quadrature gives extraction matrices
$\mathsf F^{\rightarrow}_{x_p},\mathsf F^{\leftarrow}_{x_p}\in\CC^{K\times n}$.

\subsection{Lead cells: scattering matrix and decaying subspace}\label{sec:lead}
In the right cell $\mathcal C_{0,+}$ with ports $x_{\mathrm L}=x_+$,
$x_{\mathrm R}=x_++L_+$, the incident field
$u_{\mathrm b}=\sum_m(a_{\mathrm L,m}r_m^{x_{\mathrm L}}+b_{\mathrm R,m}l_m^{x_{\mathrm R}})$
determines $\eta$ through $\mathscr J_k[\eta;u_{\mathrm b}]=0$, and the outgoing
amplitudes are $b_{\mathrm L}=\mathsf E(L_+)b_{\mathrm R}+\mathsf F^{\leftarrow}_{x_{\mathrm L}}\eta$,
$a_{\mathrm R}=\mathsf E(L_+)a_{\mathrm L}+\mathsf F^{\rightarrow}_{x_{\mathrm R}}\eta$.
This defines the cell scattering matrix
$[b_{\mathrm L};a_{\mathrm R}]=[R_{\mathrm L},T_{\mathrm{RL}};T_{\mathrm{LR}},R_{\mathrm R}]
[a_{\mathrm L};b_{\mathrm R}]$, or equivalently, for the port states
$z_{\mathrm L}=(a_{\mathrm L};b_{\mathrm L})$ and $z_{\mathrm R}=(a_{\mathrm R};b_{\mathrm R})$,
\begin{equation}\label{eq:pencil}
  \mathsf E_{\mathrm L}z_{\mathrm L}=\mathsf E_{\mathrm R}z_{\mathrm R},\qquad
  \mathsf E_{\mathrm L}=\begin{bmatrix}-R_{\mathrm L}&I\\T_{\mathrm{LR}}&0\end{bmatrix},\quad
  \mathsf E_{\mathrm R}=\begin{bmatrix}0&T_{\mathrm{RL}}\\I&-R_{\mathrm R}\end{bmatrix}.
\end{equation}
The left cell, with its own rod, contrast and width $L_-$, is treated in the
same way. Floquet solutions satisfy $\mathsf E_{\mathrm L}z=\zeta\mathsf E_{\mathrm R}z$
with multiplier $\zeta=\ee^{\ii\theta L_+}$. In a gap the continuous
multipliers pair as $\zeta\leftrightarrow1/\bar\zeta$; the pairing
$\zeta\leftrightarrow1/\zeta$ of \cite{JolyLiFliss2006} corresponds to
replacing $\beta$ by $-\beta$ and agrees with the fixed-$\beta$ problem only
under $y\mapsto-y$ symmetry. Neither implies that the discrete pencil has
exactly $K$ multipliers inside the unit circle, so this is checked.

An ordered QZ decomposition $\mathsf Q\mathsf E_{\mathrm L}\mathsf Z=\mathsf U_{\mathrm L}$,
$\mathsf Q\mathsf E_{\mathrm R}\mathsf Z=\mathsf U_{\mathrm R}$ with the $K$
eigenvalues $|\zeta|<1$ first gives the leading columns $\mathsf Z_+$ and
\begin{equation}\label{eq:propagation}
  \mathsf P_+=\mathsf U_{\mathrm R,11}^{-1}\mathsf U_{\mathrm L,11},\qquad
  \mathsf E_{\mathrm L}\mathsf Z_+=\mathsf E_{\mathrm R}\mathsf Z_+\mathsf P_+ .
\end{equation}
Hence the port states $z_j=\mathsf Z_+\mathsf P_+^jc$ of consecutive cells are
compatible and decay geometrically: the discrete analogue of the
finite-energy solutions in the right half guide. On the left, the $K$ eigenvalues $1/\zeta$ inside the unit circle of the pencil
$(\mathsf E_{\mathrm R},\mathsf E_{\mathrm L})$ give $\mathsf Z_-$ and $\mathsf P_-$ with
$\mathsf E_{\mathrm R}\mathsf Z_-=\mathsf E_{\mathrm L}\mathsf Z_-\mathsf P_-$.
Using the whole Schur subspace keeps multiple and nonnormal multipliers; a
change of basis $\mathsf Z_\pm\mapsto\mathsf Z_\pm W$ only conjugates
$\mathsf P_\pm$. If the pencil is singular, the cluster dimension is not $K$,
or the cluster is not separated from the unit circle, the discretization at
that $k$ is rejected. The discrete Cauchy data spaces are
$\mathcal G^h_\pm(k)=\operatorname{range}\mathsf C_\pm\mathsf Z_\pm$.

\subsection{Center cell}\label{sec:center}
In $\mathcal C_0$ we use \eqref{eq:cell-rep} with $\Omega=\Omega_0$ (no layer
potentials if $\Omega_0=\varnothing$) and
\[
  u_{\mathrm b}=u_{\mathrm{hom}}=\sum_m\bigl(a^-_{0,m}r_m^{x_-}+b^+_{0,m}l_m^{x_+}\bigr),
\]
the part of the field entering through the two ports. It may be nonzero even
if $\Omega_0=\varnothing$, so quasi-periodic layer potentials alone do not
represent the center field. The ansatz follows the complementary-field
representation of \cite[App.~A]{BarnettGreengard2010} with the incident port
terms added. We use it only as an ansatz; neither completeness nor uniqueness of
the representation enters the enclosure paper~\companion.

After truncation, the center unknown is $\xi=(\eta;a^-_0;b^+_0)\in\CC^{n+2K}$
($n=0$ for an empty center), and the discrete transmission equations read
$J(k)\xi=0$. The center field has the port amplitudes
\[
  \mathsf Y_-\xi=\begin{bmatrix}a^-_0\\ \mathsf F^{\leftarrow}_{x_-}\eta+\mathsf E(L_0)b^+_0\end{bmatrix},
  \qquad
  \mathsf Y_+\xi=\begin{bmatrix}\mathsf F^{\rightarrow}_{x_+}\eta+\mathsf E(L_0)a^-_0\\ b^+_0\end{bmatrix}.
\]

\subsection{Discrete eigenvalue problem and candidates}\label{sec:nep}
A candidate is a center solution whose port states lie in the decaying
subspaces: we ask for $\xi\ne0$ and $c_\pm\in\CC^K$ with
\begin{equation}\label{eq:nep}
  J(k)\xi=0,\qquad \mathsf Y_\pm\xi=\mathsf Z_\pm c_\pm .
\end{equation}
For root finding we need a scalar function that is analytic in $k$ and
independent of the basis of the stable subspaces; the determinant of the full
system \eqref{eq:nep} is neither. We therefore eliminate the center densities,
treat the center as a scattering cell with incident port amplitudes $\mathsf a\in\CC^{2K}$,
and let $\mathsf T_0(k)$, $\mathsf N_0(k)$ map $\mathsf a$ to the value and
normal-derivative coefficients of the total field at both ports. With the value
and normal-derivative blocks $\mathsf C_\pm\mathsf Z_\pm=[\mathsf D_\pm;\mathsf N_\pm]$
of the stable subspaces, the discrete DtN matrices are
$\Lambda_\pm=\mathsf N_\pm\mathsf D_\pm^{-1}$, and
\begin{equation}\label{eq:port-matrix}
  \mathsf A(k)=\mathsf N_0(k)-\diag\bigl(\Lambda_-(k),\Lambda_+(k)\bigr)\mathsf T_0(k)
  \in\CC^{2K\times2K}.
\end{equation}
On a complex neighborhood of a real interval free of Wood points, if both pencils
are regular with the stable cluster separated from the rest, and the matrices
used in the elimination are invertible, then $\det\mathsf A$ is analytic except
for poles where $\mathsf D_\pm$ or those matrices are singular. A singular
$\mathsf D_\pm$ is the discrete counterpart of a Dirichlet exceptional value
\cite[Thm.~4.1]{Fliss2013}, at which \eqref{eq:nep} remains valid. These conditions are monitored only at
interpolation nodes; they are empirical checks of the search, on which the final
interval does not depend.

\emph{Root search.} On each piece of the gap we interpolate $f(k)=\det\mathsf A(k)$
(computed by LU, as log-modulus and phase) at nested Chebyshev--Lobatto nodes of
degree $4,8,\dots,64$ until the trailing coefficients fall below $10^{-10}$ and two
off-grid checks agree to $10^{-9}$, bisecting otherwise; roots of the interpolant
come from the colleague matrix and are polished by Newton steps on the
interpolant \cite{Boyd2002,ZhaoBarnett2015}. A root is accepted if the scaled
singular-value ratio $s(\widehat k)=\sigma_{\min}/\sigma_{\max}$ of
$W_r\mathsf A(\widehat k)W_c$ is below $10^{-11}$; failures fall back to bisection
on $s(k)$. The candidate output is the right singular vector of the smallest
singular value, which determines $\mathsf a$, the center densities, and the lead
coordinates $c_\pm$. Together with $\widehat k$, $\mathsf P_\pm$, $\mathsf Z_\pm$
and the discretization parameters, these are the only inputs of
the enclosure paper~\companion. A small $s(\widehat k)$ says that a finite system is nearly
singular, not how far $\widehat k^2$ is from $\sigma(A(\beta))$.
